\chapter{Esperimenti, eventi e insiemi}

Chi lavora con l'informazione lavora con l'incertezza. Un bit trasmesso su un canale rumoroso può arrivare sbagliato,
un pacchetto può trovare la coda del router piena, una misura è sempre sporcata da un po' di rumore. Non sappiamo
\textbf{che cosa} succederà, ma possiamo dire \textbf{quanto è facile} che succeda: è il mestiere della
\term{teoria della probabilità}.

Prima di calcolare una qualsiasi probabilità serve però un linguaggio preciso per dire \emph{che cosa può succedere}.
Questo capitolo lo costruisce, e lo costruisce con uno strumento che probabilmente si conosce già: gli
\textbf{insiemi}. Alla fine del capitolo «si verificano $A$ e $B$» diventerà $A \cap B$, «non si verifica $A$»
diventerà $\compl{A}$, e potremo fare i conti sugli eventi con la stessa sicurezza con cui si fanno sui numeri.

% ============================================================
\section{Esperimenti aleatori e spazio campione}
% ============================================================

Tutto comincia da un esperimento il cui esito non si può prevedere.

\dfn{Esperimento aleatorio}{
  Un \term{esperimento aleatorio} è un'operazione, o un insieme di operazioni, il cui esito è \textbf{uno fra tanti
  risultati possibili}, nessuno dei quali si può prevedere con certezza prima di eseguirla.
}

La parola viene dal latino \emph{alea}, «dado» (è l'\emph{alea} di «\emph{alea iacta est}», il dado è tratto): lanciare
un dado è l'esperimento aleatorio per eccellenza. Ma sono esperimenti aleatori anche lanciare una moneta, estrarre una
carta da un mazzo, contare i pacchetti in coda in un router in un certo istante, misurare il rumore su un cavo.
La prima cosa da fare, davanti a un esperimento, è elencare \textbf{tutto} quello che può succedere.

\dfn{Spazio campione}{
  Lo \term{spazio campione} (o spazio dei campioni) $\Omega$ è l'insieme di \textbf{tutti i risultati possibili} di un
  esperimento aleatorio. Non deve essere per forza un insieme di numeri: per una moneta è $\{\text{Testa},
  \text{Croce}\}$.
}

$\Omega$ è la lettera greca \emph{omega} maiuscola; la minuscola $\omega$ indica invece \textbf{un singolo risultato},
cioè un elemento di $\Omega$. Visto che tutto il discorso si regge sugli insiemi, conviene rinfrescare subito le
parole di base.

\begin{richiamo}{gli insiemi, le parole di base}
Un \term{insieme} è una collezione di oggetti, detti \term{elementi}. Lo si può descrivere in due modi:
\begin{itemize}
  \item \textbf{elencando} gli elementi fra parentesi graffe: $A = \{1, 3, 5\}$;
  \item dando una \textbf{proprietà} che li caratterizza: $A = \{x \in \N : x \text{ è dispari e } x < 6\}$, che si legge
        «l'insieme degli $x$ naturali \emph{tali che} $x$ è dispari e minore di 6». I due punti si leggono «tali che».
\end{itemize}
I simboli che useremo continuamente sono questi.
\begin{center}
\renewcommand{\arraystretch}{1.25}
\begin{tabular}{@{}clP{6.6cm}@{}}
  \toprule
  \textbf{Simbolo} & \textbf{Si legge} & \textbf{Significato ed esempio} \\
  \midrule
  $x \in A$ & $x$ appartiene ad $A$ & $x$ è uno degli elementi di $A$: $3 \in \{1,3,5\}$ \\
  $x \notin A$ & $x$ non appartiene ad $A$ & $2 \notin \{1,3,5\}$ \\
  $A \subseteq B$ & $A$ è contenuto in $B$ & ogni elemento di $A$ è anche in $B$: $\{1,3\} \subseteq \{1,2,3\}$ \\
  $\emptyset$ & insieme vuoto & l'insieme che non ha nessun elemento \\
  $\card{A}$ & cardinalità di $A$ & il numero di elementi di $A$: $\card{\{1,3,5\}} = 3$ \\
  \bottomrule
\end{tabular}
\end{center}
In un insieme \textbf{conta solo chi c'è}: non conta l'ordine in cui gli elementi sono scritti e non conta se un
elemento è scritto due volte. Quindi $\{1,2\} = \{2,1\} = \{1,1,2\}$: sono tutti lo stesso insieme, con due elementi.
Due insiemi sono \textbf{uguali} quando hanno esattamente gli stessi elementi.

L'insieme vuoto è contenuto in qualsiasi insieme ($\emptyset \subseteq A$): non ha elementi, quindi non ne ha nessuno
che manchi in $A$. Anche $A \subseteq A$ è sempre vero.
\end{richiamo}

\subsection{Spazi discreti e spazi continui}

Lo spazio campione può essere di due tipi molto diversi.
\begin{itemize}
  \item È \term{discreto} se i suoi elementi si possono \textbf{mettere in fila e contare} uno dopo l'altro: il primo,
        il secondo, il terzo\ldots{} Questo accade se è \term{finito} (come le sei facce di un dado) oppure se è
        \term{infinito numerabile} (come i numeri naturali $0, 1, 2, 3, \ldots$: la fila non finisce mai, ma ogni
        elemento ha la sua posizione nella fila).
  \item È \term{continuo} se i suoi elementi sono «troppi» per essere messi in fila, come i punti di un segmento: il
        tempo, in secondi, che un pacchetto impiega ad arrivare può essere $0{,}1$, $0{,}11$, $0{,}101$\ldots{} e fra due
        valori qualsiasi ce ne sono sempre infiniti altri.
\end{itemize}

\begin{richiamo}{finito, numerabile, non numerabile}
Un insieme è \term{finito} se i suoi elementi si contano e il conteggio finisce: $\card{A}$ è un numero naturale.

Un insieme infinito è \term{numerabile} se si può fare un \textbf{elenco} che contenga tutti i suoi elementi, ognuno
in una posizione precisa. I naturali $\N_0 = \{0, 1, 2, \ldots\}$ sono numerabili per definizione. Lo sono anche i
numeri pari ($0, 2, 4, \ldots$) e persino gli interi relativi $\Z$, se li si elenca così:
\[ 0,\ 1,\ -1,\ 2,\ -2,\ 3,\ -3,\ \ldots \]
In questo elenco ogni intero compare prima o poi: il numero $-5$, per esempio, è all'undicesimo posto.

Invece i numeri reali di un intervallo, come $[0, 1]$, \textbf{non} sono numerabili: si può dimostrare (è un celebre
risultato di Georg Cantor) che qualunque elenco si provi a scrivere ne lascia fuori qualcuno.
\end{richiamo}

Per tutta la prima parte del corso, fino alle variabili aleatorie continue, lo spazio campione sarà
\textbf{discreto}: finito oppure numerabile.

\subsection{Eventi}

Spesso non interessa il risultato preciso, ma una sua caratteristica: «è uscito un numero pari», «ci sono meno di sei
pacchetti in coda». Ognuna di queste frasi individua un gruppo di risultati, cioè un sottoinsieme di $\Omega$.

\dfn{Evento}{
  Un \term{evento} è un sottoinsieme di $\Omega$, definito con una formula matematica oppure a parole. Si dice che
  l'evento $A$ \textbf{si verifica} quando il risultato $\omega$ dell'esperimento appartiene ad $A$.
}

\dfn{Evento elementare}{
  Un \term{evento elementare} è uno dei $\card{\Omega}$ risultati possibili, cioè un singolo elemento
  $\omega \in \Omega$ (o, se si preferisce vederlo come insieme, il sottoinsieme $\{\omega\}$ che contiene solo lui).
}

\begin{figure}[H]
\centering
\begin{tikzpicture}
  \draw[mspigray, line width=0.5pt, rounded corners=4pt] (-0.7,-1.0) rectangle (6.7,0.75);
  \node[mspigray, anchor=south east, font=\small] at (6.7,-1.0) {$\Omega$};
  \foreach \n in {2,4,6}
    \draw[evid!70!black, fill=evid, rounded corners=4pt] ({(\n-1)*1.2-0.45},-0.45) rectangle ({(\n-1)*1.2+0.45},0.45);
  \foreach \n in {1,...,6} {
    \node at ({(\n-1)*1.2},0) {\dado[0.62cm]{\n}};
    \node[font=\scriptsize, mspigray] at ({(\n-1)*1.2},-0.72) {$\omega = \n$};
  }
  \node[mspiA!85!black, font=\small] (lab) at (3.0,1.5) {$A = \{\text{esce un numero pari}\} = \{2, 4, 6\}$};
  \foreach \n in {2,4,6} \draw[mspiA!85!black, -{Stealth}] (lab.south) -- ({(\n-1)*1.2},0.48);
\end{tikzpicture}
\caption{Lo spazio campione del lancio di un dado ha sei eventi elementari; l'evento «esce un numero pari» è il
sottoinsieme formato dai tre risultati evidenziati.}
\end{figure}

Una cosa da tenere a mente fin da subito: un evento è determinato \textbf{solo dagli elementi che contiene}, mentre la
frase che lo descrive non è unica. Sul dado, «esce un numero pari», «esce un multiplo di 2» e «non esce un numero
dispari» sono tre frasi diverse per \textbf{lo stesso} evento $\{2, 4, 6\}$.

% ============================================================
\section{Quattro esperimenti da tenere a mente}
% ============================================================

Gli esempi che seguono torneranno per tutto il corso: conviene guardarli bene.

\subsection{Il lancio di una moneta}

Con un \textbf{solo lancio} i risultati possibili sono due:
\[ \Omega = \{\text{Testa}, \text{Croce}\} = \{T, C\}, \qquad \card{\Omega} = 2 . \]
Eventi possibili sono $A = \{T\}$ («esce testa»), $B = \{C\}$ («esce croce») e «esce testa o croce», che è
$\{T, C\} = \Omega$: un evento che si verifica sempre.

Se la moneta si lancia \textbf{due volte}, un risultato è la coppia formata dall'esito del primo lancio e da quello del
secondo:
\[ \Omega = \{TT,\ TC,\ CT,\ CC\} = \{T, C\} \times \{T, C\} = \{T, C\}^2, \qquad \card{\Omega} = 4 . \]
Il simbolo $\times$ fra due insiemi è il \textbf{prodotto cartesiano}, e merita un richiamo.

\begin{richiamo}{il prodotto cartesiano}
Il \term{prodotto cartesiano} $A \times B$ è l'insieme di tutte le \textbf{coppie ordinate} $(a, b)$ con il primo
elemento preso da $A$ e il secondo da $B$:
\[ A \times B = \{(a, b) : a \in A,\ b \in B\} . \]
«Ordinate» vuol dire che \textbf{conta l'ordine}: $(T, C)$ e $(C, T)$ sono due coppie diverse. Questa è la grande
differenza con gli insiemi, dove $\{T, C\} = \{C, T\}$. Quando un insieme si moltiplica per sé stesso si scrive una
potenza: $A \times A = A^2$, $A \times A \times A = A^3$ (le terne ordinate), e così via.

Il nome viene da Cartesio: le coppie ordinate di numeri reali, $\R \times \R = \R^2$, sono proprio i punti del piano
cartesiano, e $(2, 3)$ non è lo stesso punto di $(3, 2)$.
\end{richiamo}

\begin{figure}[H]
\centering
\begin{tikzpicture}[x=1.5cm, y=1.1cm]
  \node[font=\small\bfseries] at (-1.35,0.5) {1° lancio};
  \node[font=\small\bfseries] at (1,2.05) {2° lancio};
  \node at (0.5,1.55) {\moneta{T}}; \node at (1.5,1.55) {\moneta{C}};
  \node at (-0.45,1) {\moneta{T}}; \node at (-0.45,0) {\moneta{C}};
  \fill[evid] (1,0.5) rectangle (2,1.5);
  \fill[evid] (0,-0.5) rectangle (1,0.5);
  \draw[mspigray, shift={(0,-0.5)}, step=1] (0,0) grid (2,2);
  \node at (0.5,1) {$TT$}; \node at (1.5,1) {$TC$};
  \node at (0.5,0) {$CT$}; \node at (1.5,0) {$CC$};
  \node[font=\small, anchor=west, text width=5.4cm] at (2.25,0.5)
    {Le caselle gialle sono l'evento $A = \{\text{numero dispari di croci}\} = \{TC, CT\}$. L'evento
     $B = \{\text{nessuna croce}\}$ è la sola casella $TT$.};
\end{tikzpicture}
\caption{I quattro risultati del lancio doppio: righe per il primo lancio, colonne per il secondo.}
\end{figure}

\info{Perché \textit{TC} e \textit{CT} sono risultati diversi}{
  Lo si capisce immaginando di lanciare due monete diverse, una da un euro e una da due euro: «testa sulla moneta da
  uno e croce su quella da due» è chiaramente un'altra cosa rispetto a «croce sulla moneta da uno e testa su quella da
  due». Lanciare due volte la stessa moneta è la stessa situazione, con il primo e il secondo lancio al posto delle due
  monete. Contare $TC$ e $CT$ come un unico risultato sarebbe un errore, e nel Capitolo~3 vedremo che porta a
  probabilità sbagliate.
}

\subsection{Il lancio di un dado}

Lanciando un dado una volta:
\[ \Omega = \{1, 2, 3, 4, 5, 6\}, \qquad \card{\Omega} = 6 . \]
Ecco alcuni eventi, prima a parole e poi come insiemi:
\begin{itemize}
  \item $A$ = «il risultato è dispari» $= \{1, 3, 5\}$;
  \item $B$ = «il risultato è pari» $= \{2, 4, 6\}$;
  \item $C$ = «il risultato è dispari e non supera 4» $= \{1, 3\}$;
  \item $D$ = «il risultato è 6» $= \{6\}$, un evento elementare.
\end{itemize}

\subsection{I pacchetti in coda a un router}

Un router riceve pacchetti e li inoltra uno alla volta; quelli che arrivano mentre è occupato aspettano in una
\textbf{coda}. L'esperimento consiste nel guardare la coda in un certo istante e contare quanti pacchetti ci sono.

\begin{figure}[H]
\centering
\begin{tikzpicture}[pk/.style={draw=mspiA!85!black, fill=mspiA!20, minimum width=0.5cm, minimum height=0.36cm,
                    inner sep=0pt}]
  \draw[-{Stealth}, line width=0.8pt] (-0.8,0) -- (0,0) node[midway, above, font=\scriptsize] {arrivi};
  \draw[line width=0.7pt] (0.1,0.32) -- (3.5,0.32) -- (3.5,-0.32) -- (0.1,-0.32);
  \foreach \i in {0,...,3} \node[pk] at ({3.15-\i*0.62},0) {};
  \node[font=\scriptsize, mspigray] at (1.6,-0.55) {coda: qui ci sono 4 pacchetti};
  \node[draw, circle, fill=mspilight, minimum size=1.1cm, font=\small] (r) at (4.4,0) {router};
  \draw[-{Stealth}, line width=0.8pt] (r.east) -- ++(0.9,0) node[midway, above, font=\scriptsize] {uscita};
\end{tikzpicture}
\caption{Il numero di pacchetti in coda in un istante è il risultato di un esperimento aleatorio.}
\end{figure}

I valori possibili sono $0, 1, 2, \ldots$ e non c'è un massimo che si possa fissare a priori (o, se c'è, è così grande
che conviene ignorarlo). Lo spazio campione è allora l'insieme dei numeri naturali:
\[ \Omega = \{0, 1, 2, 3, 4, 5, 6, \ldots\} = \N_0, \qquad \card{\Omega} = \infty . \]
È uno spazio \textbf{infinito ma numerabile}, quindi discreto. Alcuni eventi:
\begin{itemize}
  \item $A$ = «in coda ci sono meno di 6 pacchetti» $= \{0, 1, 2, 3, 4, 5\}$;
  \item $B$ = «il numero di pacchetti è dispari» $= \{1, 3, 5, \ldots, 2k+1, \ldots\}$;
  \item $C$ = «il numero di pacchetti è dispari e non supera 6» $= \{1, 3, 5\}$: sono i numeri che stanno sia in $A$ sia
        in $B$, cioè (lo vedremo fra un attimo) $C = A \cap B$;
  \item $D$ = «il numero di pacchetti è pari oppure non supera 4» $= \{0, 1, 2, 3, 4, 6, 8, 10, \ldots, 2k, \ldots\}$.
\end{itemize}

\begin{figure}[H]
\centering
\begin{tikzpicture}[x=0.92cm]
  \draw[-{Stealth}, line width=0.6pt] (-0.4,0) -- (12.7,0);
  \foreach \n in {0,...,12} {
    \draw (\n,0.08) -- (\n,-0.08);
    \node[font=\scriptsize, below] at (\n,-0.14) {\n};
  }
  \node[font=\scriptsize, mspigray] at (12.45,0.25) {$\ldots$};
  % B = dispari, C = A cap B cerchiati
  \foreach \n in {1,3,...,11} \fill[mspiB] (\n,0) circle (0.11);
  \foreach \n in {0,2,...,12} \draw[fill=white] (\n,0) circle (0.11);
  \foreach \n in {1,3,5} \draw[mspiB!85!black, line width=0.8pt] (\n,0) circle (0.22);
  % A = meno di 6
  \draw[decorate, decoration={brace, amplitude=5pt}, mspiA!85!black, line width=0.8pt] (-0.2,0.35) -- (5.2,0.35)
    node[midway, above=5pt, font=\small] {$A = \{0, 1, \ldots, 5\}$};
  \node[font=\small, anchor=west, align=left] at (6.2,0.85)
    {{\color{mspiB}$\bullet$} numeri dispari: l'evento $B$\\
     {\color{mspiB!85!black}$\bigcirc$} cerchiati: $C = A \cap B = \{1, 3, 5\}$};
\end{tikzpicture}
\caption{Gli eventi del router sulla retta dei numeri naturali.}
\end{figure}

% ============================================================
\section{Operazioni fra eventi}
% ============================================================

Gli eventi sono insiemi, quindi si combinano con le operazioni degli insiemi. Ognuna ha un significato preciso a
parole, ed è il significato che conta. Nelle definizioni $A$ e $B$ sono due sottoinsiemi di $\Omega$ (più in
generale $\{A_i\}_{i=1}^{M}$ indica $M$ sottoinsiemi $A_1, A_2, \ldots, A_M$ di $\Omega$).

\dfn{Unione}{
  L'\term{unione} $A \cup B$ è l'insieme degli elementi che stanno in $A$ \textbf{oppure} in $B$ (o in entrambi: gli
  elementi comuni si prendono una volta sola). L'evento $A \cup B$ si verifica quando si verifica \textbf{almeno uno}
  fra $A$ e $B$.
}

\dfn{Intersezione}{
  L'\term{intersezione} $A \cap B$ è l'insieme di tutti e soli gli elementi comuni ad $A$ \textbf{e} a $B$. L'evento
  $A \cap B$ si verifica quando si verificano \textbf{entrambi}.
}

\dfn{Complementare}{
  Il \term{complementare} di $A$ (in $\Omega$) è l'insieme $\compl{A}$ degli elementi di $\Omega$ che \textbf{non}
  appartengono ad $A$. L'evento $\compl{A}$ si verifica quando $A$ \textbf{non} si verifica.
}

\dfn{Differenza}{
  La \term{differenza} $A \setminus B$ («$A$ meno $B$») è l'insieme degli elementi di $A$ che non appartengono a $B$: da
  $A$ si tolgono gli elementi in comune con $B$. L'evento $A \setminus B$ si verifica quando si verifica $A$
  \textbf{ma non} $B$.
}

\begin{figure}[H]
\centering
\venn{\fill[evid] \vennA \vennB;}{$A \cup B$\\ almeno uno dei due}\hfill
\venn{\begin{scope}\clip \vennA; \fill[evid] \vennB; \end{scope}}{$A \cap B$\\ tutti e due}\hfill
\venn{\fill[evid, even odd rule] \vennOmega \vennA;}{$\compl{A}$\\ non $A$}\hfill
\venn{\begin{scope}\clip \vennA; \fill[evid, even odd rule] \vennOmega \vennB; \end{scope}}{$A \setminus B$\\
  $A$ ma non $B$}
\caption{Le quattro operazioni con i \textbf{diagrammi di Venn}: il rettangolo è $\Omega$, i cerchi sono gli eventi,
la parte gialla è il risultato.}
\label{fig:operazioni}
\end{figure}

Dalla definizione di complementare seguono alcune proprietà che si usano di continuo:
\[
  \compl{\Omega} = \emptyset, \qquad \compl{\emptyset} = \Omega, \qquad \compl{\compl{A}} = A, \qquad
  A \cup \compl{A} = \Omega, \qquad A \cap \compl{A} = \emptyset .
\]
Ognuna si legge facilmente a parole: fuori da $\Omega$ non c'è niente; ciò che non sta nel «non $A$» sta in $A$; ogni
elemento sta in $A$ oppure fuori da $A$, e non può stare in entrambi.

\begin{esempio}{Una differenza}
Siano $A_1 = \{1, 2, 3, 4, 5, 6, 7, 8\}$ e $A_2 = \{1, 4, 6\}$, dentro lo spazio $\Omega = \N$. Il complementare di
$A_2$ è l'insieme di tutti i naturali diversi da $1$, $4$ e $6$. Togliendo da $A_1$ gli elementi che ha in comune con
$A_2$ resta
\[ A_1 \setminus A_2 = \{2, 3, 5, 7, 8\} . \]
Si noti che è anche l'insieme degli elementi di $A_1$ che stanno nel complementare di $A_2$: è $A_1 \cap
\compl{A_2}$. Non è un caso, come mostra la proprietà seguente.
\end{esempio}

\teorema[differenza]{La differenza come intersezione}{
  Per due eventi qualsiasi $A$ e $B$ vale
  \[ A \setminus B = A \cap \compl{B} . \]
}

\begin{dimostrazione}
\passo{Che cosa significa «uguali»} Due insiemi sono uguali se hanno gli stessi elementi. Basta quindi mostrare che un
elemento qualsiasi $x$ sta nel primo insieme se e solo se sta nel secondo.
\passo{Si traduce la definizione} $x \in A \setminus B$ vuol dire, per definizione di differenza: $x \in A$ e
$x \notin B$.
\passo{Si riconosce il complementare} «$x \notin B$» è proprio la definizione di $x \in \compl{B}$. Quindi la frase
diventa: $x \in A$ e $x \in \compl{B}$.
\passo{Si riconosce l'intersezione} «$x$ sta in $A$ e in $\compl{B}$» è la definizione di $x \in A \cap \compl{B}$.
Tutti i passaggi si possono percorrere anche all'indietro, quindi i due insiemi hanno gli stessi elementi.
\end{dimostrazione}

% ============================================================
\section{Le leggi dell'algebra degli eventi}
% ============================================================

Con i numeri sappiamo che $a \cdot (b + c) = a \cdot b + a \cdot c$ e che $(a + b) + c = a + (b + c)$. Anche le
operazioni sugli eventi obbediscono a delle leggi, e queste leggi permettono di trasformare un evento complicato in uno
più semplice da calcolare. Per dimostrarle useremo sempre la stessa tecnica, quella appena vista: tradurre ogni
operazione nelle parole «e», «o», «non».

\begin{richiamo}{«e», «o», «non» e come si dimostra che due insiemi sono uguali}
Le operazioni sugli insiemi corrispondono esattamente ai connettivi della logica:
\begin{center}
\renewcommand{\arraystretch}{1.25}
\begin{tabular}{@{}lll@{}}
  \toprule
  \textbf{Insiemi} & \textbf{Logica} & \textbf{Vale quando\ldots} \\
  \midrule
  $x \in A \cap B$ & $x \in A$ \textbf{e} $x \in B$ & sono vere tutte e due le affermazioni \\
  $x \in A \cup B$ & $x \in A$ \textbf{o} $x \in B$ & è vera almeno una delle due \\
  $x \in \compl{A}$ & \textbf{non} $x \in A$ & l'affermazione $x \in A$ è falsa \\
  \bottomrule
\end{tabular}
\end{center}
La «o» della matematica è \textbf{inclusiva}: «$P$ o $Q$» è vera anche quando sono vere tutte e due (come nell'unione,
che contiene anche gli elementi comuni). Non è la «o» esclusiva di «o mangi la minestra o salti dalla finestra».

Per dimostrare che due insiemi $X$ e $Y$ sono uguali si mostra che, per un elemento $x$ qualsiasi,
\[ x \in X \iff x \in Y , \]
dove il simbolo $\iff$ si legge «se e solo se» e significa che le due affermazioni sono o entrambe vere o entrambe
false. In pratica si scrive una \textbf{catena} di affermazioni equivalenti che parte da $x \in X$ e arriva a
$x \in Y$, giustificando ogni anello.
\end{richiamo}

\subsection{Le leggi di De Morgan}

\teorema[demorgan]{Leggi di De Morgan}{
  Per due eventi qualsiasi $A_1$ e $A_2$:
  \[
    \compl{A_1 \cup A_2} = \compl{A_1} \cap \compl{A_2}
    \qquad\qquad
    \compl{A_1 \cap A_2} = \compl{A_1} \cup \compl{A_2} .
  \]
  Il complementare di un'unione è l'intersezione dei complementari, e il complementare di un'intersezione è l'unione
  dei complementari.
}

A parole la prima legge dice una cosa ovvia: «non (piove o nevica)» è la stessa cosa di «non piove e non nevica», cioè
«né piove né nevica». La seconda: «non (sono stanco e ho fame)» vuol dire che almeno una delle due cose manca, cioè
«non sono stanco o non ho fame».

\begin{dimostrazione}
\passo{Prima legge, si parte dal primo membro} Sia $x$ un elemento qualsiasi di $\Omega$:
\[ x \in \compl{A_1 \cup A_2} \iff x \notin A_1 \cup A_2 \]
per definizione di complementare.
\passo{Si nega la «o»} $x$ non sta nell'unione quando non è vero che sta in almeno uno dei due insiemi, cioè quando
\textbf{non sta in nessuno} dei due:
\[ x \notin A_1 \cup A_2 \iff x \notin A_1 \ \text{e}\ x \notin A_2 . \]
\passo{Si arriva al secondo membro} Per definizione di complementare e poi di intersezione:
\[ x \notin A_1 \ \text{e}\ x \notin A_2 \iff x \in \compl{A_1} \ \text{e}\ x \in \compl{A_2}
   \iff x \in \compl{A_1} \cap \compl{A_2} . \]
La catena è fatta di equivalenze, quindi i due insiemi coincidono.
\passo{La seconda legge si ricava dalla prima} La prima legge vale per due eventi qualsiasi, quindi anche per i due
eventi $\compl{A_1}$ e $\compl{A_2}$:
\[ \compl{\compl{A_1} \cup \compl{A_2}} = \compl{\compl{A_1}} \cap \compl{\compl{A_2}} = A_1 \cap A_2 , \]
perché il complementare del complementare è l'insieme di partenza.
\passo{Si fa il complementare di entrambi i membri} Due insiemi uguali hanno complementari uguali:
\[ \compl{\compl{\compl{A_1} \cup \compl{A_2}}} = \compl{A_1 \cap A_2}
   \quad\Longrightarrow\quad \compl{A_1} \cup \compl{A_2} = \compl{A_1 \cap A_2} , \]
che è la seconda legge.
\end{dimostrazione}

\begin{figure}[H]
\centering
\venn{\fill[evid, even odd rule] \vennOmega \vennA;}{$\compl{A}$}\hfill
\venn{\fill[evid, even odd rule] \vennOmega \vennB;}{$\compl{B}$}\hfill
\venn{\fill[evid, even odd rule] \vennOmega \vennA; \fill[white] \vennB;}{$\compl{A} \cap \compl{B}$\\
  (giallo in entrambe)}\hfill
\venn{\fill[evid, even odd rule] \vennOmega \vennA; \fill[white] \vennB;}{$\compl{A \cup B}$\\
  (fuori dall'unione)}
\caption{La prima legge di De Morgan vista con i diagrammi: la parte gialla in comune ai primi due disegni è proprio
ciò che resta fuori dall'unione.}
\end{figure}

\subsection{Le proprietà associativa e distributiva}

\teorema[associativa]{Proprietà associativa}{
  Per tre eventi qualsiasi $A_1$, $A_2$, $A_3$:
  \[
    (A_1 \cup A_2) \cup A_3 = A_1 \cup (A_2 \cup A_3)
    \qquad\qquad
    (A_1 \cap A_2) \cap A_3 = A_1 \cap (A_2 \cap A_3) .
  \]
}

\begin{dimostrazione}
\passo{Unione} Un elemento sta in $(A_1 \cup A_2) \cup A_3$ se sta in $A_1 \cup A_2$ oppure in $A_3$, cioè se sta in
\textbf{almeno uno} dei tre insiemi. Lo stesso ragionamento, fatto sul secondo membro, porta alla stessa frase.
\passo{Intersezione} Un elemento sta in $(A_1 \cap A_2) \cap A_3$ se sta in $A_1 \cap A_2$ e in $A_3$, cioè se sta in
\textbf{tutti e tre} gli insiemi; e di nuovo il secondo membro dice lo stesso.
\end{dimostrazione}

Grazie a questa proprietà le parentesi non servono più: si scrive semplicemente $A_1 \cup A_2 \cup A_3$ («almeno uno
dei tre») e $A_1 \cap A_2 \cap A_3$ («tutti e tre»). Per molti insiemi si usa una notazione compatta, la stessa idea
della sommatoria che vedremo nel prossimo capitolo.

\begin{richiamo}{unione e intersezione di molti insiemi}
Dati $M$ insiemi $A_1, A_2, \ldots, A_M$ si scrive
\[
  \bigcup_{i=1}^{M} A_i = A_1 \cup A_2 \cup \cdots \cup A_M ,
  \qquad\qquad
  \bigcap_{i=1}^{M} A_i = A_1 \cap A_2 \cap \cdots \cap A_M .
\]
Si legge «unione (intersezione) per $i$ che va da 1 a $M$ degli $A_i$». La lettera $i$ è un \term{indice}: prende uno
dopo l'altro i valori $1, 2, \ldots, M$, e per ciascuno si prende l'insieme $A_i$. Il primo insieme contiene gli
elementi che stanno in \textbf{almeno uno} degli $A_i$; il secondo quelli che stanno in \textbf{tutti}. L'indice può
partire da un altro numero: $\bigcap_{i=2}^{M} A_i = A_2 \cap A_3 \cap \cdots \cap A_M$.
\end{richiamo}

\teorema[distributiva]{Proprietà distributive}{
  Per eventi qualsiasi $A_1, A_2, \ldots, A_M$ valgono la distributiva dell'unione rispetto all'intersezione e quella
  dell'intersezione rispetto all'unione:
  \[
    A_1 \cup \left( \bigcap_{i=2}^{M} A_i \right) = \bigcap_{i=2}^{M} \left( A_1 \cup A_i \right)
    \qquad\qquad
    A_1 \cap \left( \bigcup_{i=2}^{M} A_i \right) = \bigcup_{i=2}^{M} \left( A_1 \cap A_i \right) .
  \]
  Con tre soli eventi $A$, $B$, $C$ si leggono meglio:
  \[
    A \cup (B \cap C) = (A \cup B) \cap (A \cup C)
    \qquad\qquad
    A \cap (B \cup C) = (A \cap B) \cup (A \cap C) .
  \]
}

\begin{dimostrazione}
\passo{Seconda legge, primo membro} $x \in A_1 \cap \left(\bigcup_{i=2}^{M} A_i\right)$ vuol dire che $x \in A_1$
\textbf{e} che $x$ sta in almeno uno fra $A_2, \ldots, A_M$: diciamo in $A_k$.
\passo{Seconda legge, secondo membro} Allora $x$ sta sia in $A_1$ sia in $A_k$, cioè in $A_1 \cap A_k$: sta in almeno
una delle intersezioni $A_1 \cap A_i$, quindi nella loro unione. Viceversa, se $x$ sta in una delle $A_1 \cap A_k$,
sta in $A_1$ e in $A_k$, quindi in $A_1$ e nell'unione degli $A_i$. I due insiemi coincidono.
\passo{Prima legge, primo membro} $x \in A_1 \cup \left(\bigcap_{i=2}^{M} A_i\right)$ vuol dire che $x \in A_1$
\textbf{oppure} che $x$ sta in tutti gli $A_2, \ldots, A_M$.
\passo{Prima legge, secondo membro} Se $x \in A_1$, allora $x$ sta in ogni $A_1 \cup A_i$ (che contiene $A_1$); se
invece $x$ sta in tutti gli $A_i$, sta di nuovo in ogni $A_1 \cup A_i$. In entrambi i casi $x$ sta nell'intersezione
di tutti gli $A_1 \cup A_i$.
\passo{Prima legge, il viceversa} Sia ora $x$ in tutti gli $A_1 \cup A_i$. Se $x \in A_1$ abbiamo finito. Se
$x \notin A_1$, allora da $x \in A_1 \cup A_i$ segue che $x \in A_i$, e questo per ogni $i$ da 2 a $M$: $x$ sta
nell'intersezione degli $A_i$. Anche qui i due insiemi coincidono.
\end{dimostrazione}

\begin{figure}[H]
\centering
\venntre{\begin{scope}\clip \vennTA; \fill[evid] \vennTB \vennTC; \end{scope}}{$A \cap (B \cup C)$}\hspace{1.2cm}
\venntre{\begin{scope}\clip \vennTA; \fill[evid] \vennTB; \end{scope}
         \begin{scope}\clip \vennTA; \fill[evid] \vennTC; \end{scope}}{$(A \cap B) \cup (A \cap C)$}
\caption{La seconda proprietà distributiva: in entrambi i casi si colora la parte di $A$ che sta anche in $B$ o in $C$.}
\end{figure}

\info{Insiemi e numeri: somiglianze e differenze}{
  Se si sostituisce l'unione con la somma e l'intersezione con il prodotto, la seconda legge distributiva diventa
  quella che si usa da sempre con i numeri: $a \cdot (b + c) = a \cdot b + a \cdot c$. La \textbf{prima} invece
  diventerebbe $a + b \cdot c = (a + b) \cdot (a + c)$, che con i numeri è \textbf{falsa}: con $a = 1$, $b = 2$,
  $c = 3$ il primo membro vale $1 + 6 = 7$ e il secondo $3 \cdot 4 = 12$. Con gli insiemi l'unione e l'intersezione
  sono perfettamente simmetriche, e ogni legge vale anche scambiando i due simboli.
}

% ============================================================
\section{Un po' di nomenclatura}
% ============================================================

Alcuni eventi e alcune relazioni fra eventi hanno un nome proprio.

\dfn{Evento certo, evento impossibile}{
  $\Omega$ è detto \term{evento certo}: contiene tutti i risultati, quindi si verifica sempre. $\emptyset$ è detto
  \term{evento impossibile}: non contiene nessun risultato, quindi non si verifica mai.
}

\dfn{Eventi complementari, incompatibili, implicazione}{
  \begin{itemize}
    \item $A$ e $\compl{A}$ sono detti \term{eventi complementari}.
    \item Due eventi $A$ e $B$ con $A \cap B = \emptyset$ sono detti \term{incompatibili} (o \term{mutuamente
          esclusivi}): non si possono verificare insieme.
    \item Se $A \subseteq B$ si dice che \term{$A$ implica $B$}: ogni volta che si verifica $A$, si verifica
          necessariamente anche $B$.
  \end{itemize}
}

\begin{figure}[H]
\centering
\begin{tikzpicture}
  \draw[mspigray, line width=0.5pt, rounded corners=2pt] (-1.6,-1.0) rectangle (1.6,1.0);
  \draw[mspiA!85!black, line width=0.8pt, fill=mspiA!12] (-0.75,0) circle (0.6);
  \draw[mspiB!85!black, line width=0.8pt, fill=mspiB!12] (0.75,0) circle (0.6);
  \node[mspiA!85!black] at (-0.75,0) {$A$}; \node[mspiB!85!black] at (0.75,0) {$B$};
  \node[mspigray, font=\footnotesize, anchor=south east] at (1.6,-1.0) {$\Omega$};
  \node[font=\small, anchor=north] at (0,-1.08) {incompatibili: $A \cap B = \emptyset$};
\end{tikzpicture}\hspace{1.5cm}
\begin{tikzpicture}
  \draw[mspigray, line width=0.5pt, rounded corners=2pt] (-1.6,-1.0) rectangle (1.6,1.0);
  \draw[mspiB!85!black, line width=0.8pt, fill=mspiB!12] (0,0) circle (0.85);
  \draw[mspiA!85!black, line width=0.8pt, fill=mspiA!15] (-0.25,-0.1) circle (0.45);
  \node[mspiA!85!black] at (-0.25,-0.1) {$A$}; \node[mspiB!85!black] at (0.45,0.45) {$B$};
  \node[mspigray, font=\footnotesize, anchor=south east] at (1.6,-1.0) {$\Omega$};
  \node[font=\small, anchor=north] at (0,-1.08) {$A$ implica $B$: $A \subseteq B$};
\end{tikzpicture}
\caption{Due eventi incompatibili non hanno risultati in comune; se $A \subseteq B$, chi sta in $A$ sta anche in $B$.}
\end{figure}

\begin{esempio}{Testa e croce sono complementari}
Nel lancio singolo di una moneta, $T$ e $C$ sono eventi complementari ($\compl{T} = C$), quindi anche incompatibili:
\[ T \cap C = \compl{C} \cap C = \emptyset . \]
Applicando le leggi di De Morgan si ritrovano le proprietà di $\Omega$ e di $\emptyset$:
\[
  \compl{T \cap C} = \compl{T} \cup \compl{C} = C \cup T = \Omega = \compl{\emptyset},
  \qquad
  \compl{T \cup C} = \compl{T} \cap \compl{C} = C \cap T = \emptyset = \compl{\Omega} .
\]
\end{esempio}

\begin{esempio}{Lo stesso evento, tante frasi}
Nel lancio doppio della moneta, l'evento $D$ = «esce almeno una croce» è incompatibile con $\{TT\}$ e si può scrivere
in molti modi:
\[
  D = \{CC, TC, CT\} = \compl{\{TT\}} = \Omega \setminus \{TT\} = \{CC\} \cup \{TC\} \cup \{CT\} .
\]
Matematicamente l'evento è individuato in modo unico dai suoi elementi; a parole lo si può descrivere con frasi
diverse: «almeno una croce», «non due teste», «non è uscito $TT$».
\end{esempio}

Per finire, ecco il «vocabolario» che traduce le frasi sugli eventi nelle operazioni sugli insiemi. Conviene impararlo
bene: negli esercizi il primo passo è sempre tradurre il testo in questa lingua.

\begin{table}[H]
\centering
\renewcommand{\arraystretch}{1.3}
\begin{tabular}{@{}ll@{}}
  \toprule
  \textbf{A parole} & \textbf{Con gli insiemi} \\
  \midrule
  si verifica $A$ & il risultato $\omega$ sta in $A$: $\omega \in A$ \\
  si verifica $A$ \textbf{oppure} $B$ (almeno uno) & $A \cup B$ \\
  si verificano $A$ \textbf{e} $B$ (entrambi) & $A \cap B$ \\
  \textbf{non} si verifica $A$ & $\compl{A}$ \\
  si verifica $A$ \textbf{ma non} $B$ & $A \setminus B = A \cap \compl{B}$ \\
  \textbf{né} $A$ \textbf{né} $B$ & $\compl{A} \cap \compl{B} = \compl{A \cup B}$ \\
  $A$ e $B$ non possono verificarsi insieme & $A \cap B = \emptyset$ \\
  se si verifica $A$, si verifica anche $B$ & $A \subseteq B$ \\
  qualcosa succede sempre / non succede mai & $\Omega$ / $\emptyset$ \\
  \bottomrule
\end{tabular}
\caption{Dal linguaggio degli eventi a quello degli insiemi.}
\end{table}

\gbox{In sintesi}{mspiA!80!black}{
\begin{itemize}
  \item Un \textbf{esperimento aleatorio} ha un esito imprevedibile; lo \textbf{spazio campione} $\Omega$ raccoglie
        tutti i risultati possibili. Per ora $\Omega$ è \textbf{discreto}: finito o numerabile.
  \item Un \textbf{evento} è un sottoinsieme di $\Omega$; si verifica quando il risultato gli appartiene. È determinato
        dai suoi elementi, non dalla frase che lo descrive.
  \item Unione («o»), intersezione («e»), complementare («non») e differenza («ma non») costruiscono nuovi eventi, e
        $A \setminus B = A \cap \compl{B}$.
  \item Leggi di \textbf{De Morgan}: $\compl{A \cup B} = \compl{A} \cap \compl{B}$ e
        $\compl{A \cap B} = \compl{A} \cup \compl{B}$. Valgono anche le proprietà associativa e le due distributive.
  \item $\Omega$ è l'evento certo, $\emptyset$ l'impossibile; $A$ e $B$ sono incompatibili se $A \cap B = \emptyset$;
        $A$ implica $B$ se $A \subseteq B$.
\end{itemize}
}
